\chapter{Experimental Details}
\label{app:construction}

This appendix specifies the population construction, exact parameter grids,
uncertainty intervals and validity rules used in the experiments.

\section{Additive construction}

For prescribed positive row probabilities \(a=(a_1,\ldots,a_r)\) and
column probabilities \(b=(b_1,\ldots,b_c)\), form the independence table
\(B=ab^{\mathsf T}\), whose cell probabilities are \(B_{ij}=a_i b_j\).
Let \(H\) specify which cells gain or lose probability and \(t\ge0\)
control the amount transferred. Define
\begin{equation}
 R(t)=ab^{\mathsf T}+tH,
 \qquad H\mathbf{1}_c=0,
 \qquad \mathbf{1}_r^{\mathsf T}H=0.
 \label{eq:additive-construction}
\end{equation}
Here \(\mathbf{1}_k\) is a vector of \(k\) ones. The zero row and column
sums of \(H\) ensure that the transfers leave both margins unchanged.
The principal first-cell direction is zero outside rows and columns 1
and 2, where it has the block
\begin{equation}
 H_{1:2,1:2}=\begin{pmatrix}1&-1\\-1&1\end{pmatrix}.
 \label{eq:first-change-block}
\end{equation}
The rare-cell direction places the same block in the last two rows and
columns. The spread direction is
\begin{equation}
 H_{ij}=u_i v_j,
 \quad u_i=-1+\frac{2(i-1)}{r-1},
 \quad v_j=-1+\frac{2(j-1)}{c-1}.
 \label{eq:appendix-spread}
\end{equation}
Both score vectors sum to zero, so every direction preserves margins.

The largest nonnegative strength before a decreasing cell reaches zero is
\begin{equation}
 t_{\max}=\min_{H_{ij}<0}\frac{B_{ij}}{-H_{ij}}.
 \label{eq:design-tmax}
\end{equation}
Evaluate the path endpoint MI\@. A target MI must lie strictly below this
endpoint, within the stated numerical tolerance, so the chosen strength
satisfies \(0\le t<t_{\max}\). This endpoint is specific to the chosen
transfer pattern, not the largest possible MI over all tables with those
margins. For a positive target,
Brent's one-dimensional root search \citep{brent1973} solves \(I\{R(t)\}=I_{\rm target}\)
on \([0,t_{\max}]\); a zero target uses \(t=0\). The final table is accepted
only if every cell is strictly positive, the sum is one, both margins agree
with their prescribed vectors, and achieved MI agrees with the target. The
normalisation and margin tolerances are \(10^{-12}\); the MI tolerance is
\(10^{-12}+10^{-8}I_{\rm target}\).

\medskip
\thesisminiheading{Equivalence to the binary sampling mechanism.}

In the binary example, both variables have category probabilities
\(a=(a_1,a_2)\). With probability \(w\), \(Y\) copies \(X\), giving
joint probabilities \(a_1,a_2\) on the diagonal and zero elsewhere.
With probability \(1-w\), \(Y\) is drawn independently, giving
\(B=aa^{\mathsf T}\). Combining these two possibilities
gives
\begin{align}
 R
 &= (1-w)
    \begin{pmatrix}a_1^2&a_1a_2\\a_1a_2&a_2^2\end{pmatrix}
    +w\begin{pmatrix}a_1&0\\0&a_2\end{pmatrix}\\
 &= \begin{pmatrix}a_1^2&a_1a_2\\a_1a_2&a_2^2\end{pmatrix}
    +wa_1a_2\begin{pmatrix}1&-1\\-1&1\end{pmatrix}.
 \label{eq:binary-copy-equivalence}
\end{align}
The second line uses \(a_1+a_2=1\), so
\(a_1-a_1^2=a_1a_2\) and \(a_2-a_2^2=a_1a_2\).
Thus the copying mechanism is exactly the first-cell additive
construction with \(t=wa_1a_2\). For the margins in Chapter~\ref{chap:design},
\(t_P\simeq0.04272945\) and \(t_Q\simeq0.04174587\) give
\(w_P\simeq0.20347356\) and \(w_Q\simeq0.26091171\), respectively.

\medskip

Row and column margins follow Table~\ref{tab:margin-profiles}, with its
formula applied separately to the two dimensions. The extreme regimes
specify their dominant probabilities directly in Table~\ref{tab:exact-grids}.

\section{Log-linear comparison}

The log-linear comparison uses the scores \(u_i,v_j\) from
Equation~\eqref{eq:appendix-spread}, but starts from positive cell weights
\begin{equation}
  K_{ij}(\lambda)=\exp(\lambda u_i v_j)
\end{equation}
instead of adding \(tH\) to the independence table \(B\). Iterative
proportional fitting \citep{deming1940} scales rows and columns of these
weights until the prescribed margins are restored. At \(\lambda=0\), this
gives the independence table; a numerical search then chooses \(\lambda\)
to reach the target MI\@. The final table passes the same positivity,
normalisation, margin and target-MI checks as the additive table. Unlike
the first-block additive construction used in the comparison, it changes
cells throughout the table. Its cell probabilities can therefore differ
from additive tables with the same margins and MI\@.

\section{Statistical experiment grids}

Unless a row of Table~\ref{tab:exact-grids} says otherwise, both populations
use the first-cell additive direction, \(I(P)=I_0\),
\(I(Q)=I_0+|\Delta|\), and the three standard margin profiles. The
symbol \(n\) means \(n_P=n_Q=n\). Only targets producing strictly positive
population tables are used.

\begin{longtable}{@{}p{0.18\textwidth}p{0.75\textwidth}@{}}
\caption{Exact statistical experiment grids; MI values are in nats.}
\label{tab:exact-grids}\\
\toprule
Family & Settings\\
\midrule
\endfirsthead
\toprule
Family & Settings\\
\midrule
\endhead
Main & Shapes \(2\times2,3\times3,5\times5,8\times8\); three profiles;
\(I_0=0.02\); \(|\Delta|\in\{0,0.001,0.002,0.005,0.01,0.02\}\);
\(n\in\{10,20,50,100,250,500,1000\}\).\\
Baseline & Main shapes and profiles; \(I_0\in\{0.0001,0.001,0.02\}\);
main difference grid; \(n\in\{20,100,1000\}\).\\
Changed cells & Shapes \(3\times3,5\times5,8\times8\);
different-skew margins; \(I_0=0.0001\);
\(|\Delta|\in\{0,0.0001,0.00025,0.0005,0.001,0.002\}\);
\(n\in\{20,100,1000\}\). First and rare blocks for all shapes;
spread pattern for \(5\times5\) and \(8\times8\).\\
Imbalance & Main shapes, profiles and difference grid;
\(I_0=0.02\). Sample pairs are
\((50,50),(50,100),(100,50),(50,250)\),
\((250,50),(50,500),(500,50),(500,500)\).
Both \(Q\)-higher and \(P\)-higher MI directions.\\
Larger effects & Shapes \(2\times2,3\times3\); three profiles;
\(I_0=0.02\); main difference grid plus \(0.05,0.1,0.2\);
\(n\in\{20,100,1000\}\).\\
Rectangular & Shapes \(2\times3,3\times5\); three profiles;
\(I_0=0.02\); main difference grid;
\(n\in\{20,100,1000\}\).\\
Construction & Shapes \(3\times3,5\times5,8\times8\); uniform and
different-skew profiles; \(I_0\in\{0.0001,0.02\}\); main difference
grid; \(n\in\{20,100,1000\}\); additive first block and ordinal
log-linear constructors.\\
Extreme skew & Shapes \(2\times2,3\times3,8\times8\); dominant margin
probability pairs \((\rho_P,\rho_Q)\in\{(0.9,0.95),(0.99,0.995),
(0.999,0.9995)\}\); \(I_0=0.00001\);
\(|\Delta|\in\{0,0.00001,0.00005,0.0001,0.0002\}\);
\(n\in\{20,100,1000,10000\}\).\\
Rare-cell stress & Shapes \(3\times3,8\times8\); dominant pair
\((0.9,0.95)\); first and rare blocks; \(I_0=0.00001\);
\(|\Delta|\in\{0,0.00001,0.000025,0.00005,0.0001\}\);
same \(n\) ladder as extreme skew.\\
Independence & Shapes \(2\times2,8\times8\); three profiles;
\(I_0=0\); main difference grid;
\(n\in\{20,100,1000\}\).\\
Convergence & Main square shapes and profiles; equal MI
\(I(P)=I(Q)\in\{0.0001,0.02\}\);
\(n\in\{1000,2500,10000,50000\}\). Additional rare-block controls
use \(3\times3,5\times5,8\times8\), different-skew margins and common
MI 0.0001.\\
Runtime & Shapes \(2\times2,3\times3,3\times5,5\times5,8\times8\);
uniform and different-skew profiles; \(I_0=0.02\);
\(|\Delta|\in\{0,0.02\}\); \(n\in\{20,100,1000\}\).
Each of the 60 regimes uses 200 pre-generated table pairs, 20 warm-up calls
and five timed passes.\\
\bottomrule
\end{longtable}

\section{Rejection rates and simulation uncertainty}
\label{sec:simulation-uncertainty}

The plotted rejection rate divides valid rejections by all 20,000 sampled
table pairs. The method-valid conditional rate instead divides by that
method's valid outputs; it is undefined if there are none. The common-valid
rate restricts the calculation to pairs for which both methods are valid.
Comparing unconditional and common-valid rates separates the effect of the
reference distribution from lost validity.

For the paired comparison, let \(D_s\) equal the Wald valid-rejection
indicator minus the Expanded Welch valid-rejection indicator for sampled
pair \(s\), matching the sign in Table~\ref{tab:null-paired-examples}.
Its average is \(\bar D=(1/20{,}000)\sum_{s=1}^{20{,}000}D_s\), with
approximate 95\% interval
\begin{equation}
 \bar D\ \pm\ 1.96
 \left\{\frac{1}{20{,}000(19{,}999)}
 \sum_{s=1}^{20{,}000}(D_s-\bar D)^2\right\}^{1/2}.
\end{equation}
Calculating \(D_s\) from each sampled pair accounts for the two methods
being applied to the same data. Individual-method
intervals use Wilson's binomial score construction \citep{wilson1927}.
Both types of interval describe Monte Carlo uncertainty for one setting,
not simultaneous uncertainty across all settings.

\section{Validity and sparse-sample diagnostics}
\label{sec:validity-rules}

Both tests require matching table shapes with at least two rows and columns,
finite nonnegative integer counts, and more than one observation in each
table. Normal Wald requires a finite corrected MI difference, a positive
finite standard error, \(\Vhat(P)+\Vhat(Q)>10^{-14}\), and a finite p-value.
The variance threshold treats values extremely close to zero as numerically
unusable.

Expanded Welch adds component-level conditions to those shared requirements.
It requires positive finite degrees of freedom for both populations and
their combination. If one component's degrees of freedom cannot be
calculated, the expanded test is invalid even when the direct combined
expression in Section~\ref{sec:zero-component} has a value. Failed
calculations contribute no rejection to the unconditional rate.

For population \(P\) with sample size \(n_P\), the minimum expected cell
count is \(\min_{ij}n_Pp_{ij}\); define \(Q\)'s analogously. Sparsity
diagnostics include empty observed rows and columns and the fraction of
zero-count cells. They help describe why a statistic might fail but are not
used to exclude settings from the study.

\section{Independent-pilot standard-error check}

The diagnostic in Section~\ref{sec:mechanism-results} uses two independent
sets of 20,000 sampled table pairs from the stated null populations. The
pilot set gives the sample standard deviation of the corrected MI
differences. The evaluation set is tested twice: once with each pair's
plug-in standard error, and once with that fixed pilot standard deviation.
Both calculations use the standard normal distribution and the two-sided
5\% threshold. The comparison of variance estimates divides the evaluation
set's average squared plug-in standard error by the squared pilot standard
deviation. This check retains the population pair and estimated differences
while changing how they are standardised.
